\makeatletter
\def\input@path{{./}{../../../preamble/}{book2026-2027/preamble/}}
\makeatother
\input{preamble.tex}
\input{graphs/highlights.tex}
% \GTenHighlightsfalse % Uncomment to disable instructional highlights.
\input{graphs/specific_highlights.tex}
\input{graphs/G10_T2_bar_graphs.tex}
% Graph markers follow the local written roles: outcomes on category labels;
% variable values on numerical ticks; bar fills and height labels follow the
% vertical quantity (blue for probabilities, the variable marker for values).
% Grouped PMFs keep probability blue; axis labels and captions identify units.
\usepackage{needspace}
\hypersetup{pdftitle={Grade 10 Term 2 Discrete Probability}}
\setlength{\emergencystretch}{2em}
\title{Grade 10, Term 2, Lesson 1, Criterion C1\\Outcomes, Support and Probability Mass Functions}
\author{}
\date{}
\begin{document}
\maketitle
\vspace{-2\baselineskip}
\textbf{C1:} Characterises discrete random variables in economic and financial contexts by identifying their support, probability mass function and possible outcomes.

\setcounter{tocdepth}{1}
\setcounter{secnumdepth}{2}
\phantomsection\label{toc}
\tableofcontents

\section{Introduction}\label{sec:t2-c1-introduction}

\SectionProblemsTableOfContents{
	\SectionProblemLink{subsec:t2-c1-given-outcomes}{Given outcomes and their probabilities} \\
	\SectionProblemLink{subsec:t2-c1-independent-outcomes}{Outcomes from independent events}
} \vspace{5mm}

All data in this activity are hypothetical teaching models. Keep exact expressions where useful and round final probabilities to four decimal places.

\PracticeHighlightLegend{G10-T2-C1}

\Needspace{12\baselineskip}
\subsection{Given outcomes and their probabilities}
\label{subsec:t2-c1-given-outcomes}

A school shop receives one payment. Its method is the \emph{complete outcome}. A \emph{random variable} assigns a number to each outcome. The table gives the outcome probabilities and maps each method to the fee $\GTenHLValue{F}$ in USD, an electronic-payment indicator $\GTenHLSecond{E}$, and a card-payment indicator $\GTenHLThird{C}$. Each indicator is 1 when its condition is met and 0 otherwise.

\begin{center}
\begin{tabular}{@{}lcccc@{}}
\toprule
\textbf{Outcome $\GTenHLOutcome{\omega}$} & $P(\GTenHLOutcome{\omega})$ & $\GTenHLValue{F}(\GTenHLOutcome{\omega})$ (USD) & $\GTenHLSecond{E}(\GTenHLOutcome{\omega})$ & $\GTenHLThird{C}(\GTenHLOutcome{\omega})$ \\
\midrule
\GTenHLOutcome{Cash}          & \GTenHLProbability{0.1} & \GTenHLValue{0} & \GTenHLSecond{0} & \GTenHLThird{0} \\
\GTenHLOutcome{Debit card}    & \GTenHLProbability{0.2} & \GTenHLValue{1} & \GTenHLSecond{1} & \GTenHLThird{1} \\
\GTenHLOutcome{Credit card}   & \GTenHLProbability{0.3} & \GTenHLValue{1} & \GTenHLSecond{1} & \GTenHLThird{1} \\
\GTenHLOutcome{Mobile wallet} & \GTenHLProbability{0.4} & \GTenHLValue{2} & \GTenHLSecond{1} & \GTenHLThird{0} \\
\bottomrule
\end{tabular}
\end{center}

The \emph{support} is the set of values with positive probability. Here it is $\GTenHLValue{\{0,1,2\}}$ for the fee in USD and $\GTenHLSecond{\{0,1\}}$ and $\GTenHLThird{\{0,1\}}$ for the unitless indicators $\GTenHLSecond{E}$ and $\GTenHLThird{C}$, respectively.

\Needspace{9\baselineskip}
A \emph{probability mass function} (PMF) gives the probability of each value. Read the complete-outcome probabilities from the table and add those of the mutually exclusive outcomes that map to the same value. In particular, a 1 USD fee and a card payment each group debit and credit together. The electronic-payment group also includes the mobile wallet.
\[
\begin{aligned}
P(\GTenHLValue{F}=\GTenHLValue{1})&=\GTenHLProbability{0.2}+\GTenHLProbability{0.3}=\GTenHLProbability{0.5},\\
P(\GTenHLSecond{E}=\GTenHLSecond{1})&=\GTenHLProbability{0.2}+\GTenHLProbability{0.3}+\GTenHLProbability{0.4}=\GTenHLProbability{0.9},\\
P(\GTenHLThird{C}=\GTenHLThird{1})&=\GTenHLProbability{0.2}+\GTenHLProbability{0.3}=\GTenHLProbability{0.5}.
\end{aligned}
\]
The fee PMF has masses $\GTenHLProbability{0.1},\GTenHLProbability{0.5},\GTenHLProbability{0.4}$ at $\GTenHLValue{0},\GTenHLValue{1},\GTenHLValue{2}$.
The indicators share a support but have different masses: $\GTenHLProbability{0.1},\GTenHLProbability{0.9}$ for $\GTenHLSecond{E}$ and $\GTenHLProbability{0.5},\GTenHLProbability{0.5}$ for $\GTenHLThird{C}$, in order from 0 to 1. All probabilities are unitless. See the \hyperref[subsec:t2-c1-introductory-general-problem-grouping-outcomes]{worked grouping problem} for the full fee calculation.

\secInicio[subsec:t2-c1-given-outcomes][sec:t2-c1-introduction]

\Needspace{12\baselineskip}
\subsection{Outcomes from independent events}
\label{subsec:t2-c1-independent-outcomes}

A school enterprise offers a 5 USD notebook to each of three visitors. For each visitor, the possible decisions are to buy (\GTenHLOutcome{B}, probability \GTenHLProbability{0.4}) or to pass (\GTenHLOutcome{N}, probability \GTenHLProbability{0.6}). Assume the visitors decide independently: one visitor's decision does not change the probabilities for the others.

Choose B or N for each visitor in order. The eight complete outcomes are
\[
\GTenHLOutcome{\Omega}=\{\GTenHLOutcome{\mathrm{NNN}},\GTenHLOutcome{\mathrm{BNN}},\GTenHLOutcome{\mathrm{NBN}},\GTenHLOutcome{\mathrm{NNB}},
\GTenHLOutcome{\mathrm{BBN}},\GTenHLOutcome{\mathrm{BNB}},\GTenHLOutcome{\mathrm{NBB}},\GTenHLOutcome{\mathrm{BBB}}\}.
\]
Multiply the three component probabilities to find the probability of one complete outcome. For example, in \GTenHLOutcome{BNN} only the first visitor buys, so
\[
P(\GTenHLOutcome{\mathrm{BNN}})=\GTenHLProbability{0.4}\times\GTenHLProbability{0.6}\times\GTenHLProbability{0.6}=\GTenHLProbability{0.144}.
\]
The product has the same factors for \GTenHLOutcome{NBN} and \GTenHLOutcome{NNB}. The products are $\GTenHLProbability{0.6}^3=\GTenHLProbability{0.216}$ for \GTenHLOutcome{NNN}, $\GTenHLProbability{0.4}^2\times\GTenHLProbability{0.6}=\GTenHLProbability{0.096}$ for each of \GTenHLOutcome{BBN}, \GTenHLOutcome{BNB}, and \GTenHLOutcome{NBB}, and $\GTenHLProbability{0.4}^3=\GTenHLProbability{0.064}$ for \GTenHLOutcome{BBB}.

\Needspace{10\baselineskip}
\begin{itemize}
	\item $\GTenHLValue{X}$ counts the buyers. Its support is $\GTenHLValue{\{0,1,2,3\}}$, measured in buyers.
	\item $\GTenHLSecond{Y}=5\GTenHLValue{X}$ is the revenue in USD, with support $\GTenHLSecond{\{0,5,10,15\}}$.
	\item $\GTenHLThird{Z}$ equals 1 if the first visitor buys and 0 otherwise. Its support is $\GTenHLThird{\{0,1\}}$ and it has no units.
\end{itemize}

Both \GTenHLOutcome{BNN} and \GTenHLOutcome{NBN} give $\GTenHLValue{X}=\GTenHLValue{1}$ and $\GTenHLSecond{Y}=\GTenHLSecond{5}$. Their indicator values differ: $\GTenHLThird{Z}(\GTenHLOutcome{\mathrm{BNN}})=\GTenHLThird{1}$ and $\GTenHLThird{Z}(\GTenHLOutcome{\mathrm{NBN}})=\GTenHLThird{0}$.

\Needspace{8\baselineskip}
To get the PMF entry for exactly one buyer, add the probabilities of the three mutually exclusive complete outcomes that give $\GTenHLValue{X}=\GTenHLValue{1}$:
\[
\begin{aligned}
P(\GTenHLValue{X}=\GTenHLValue{1})
&=P(\GTenHLOutcome{\mathrm{BNN}})+P(\GTenHLOutcome{\mathrm{NBN}})+P(\GTenHLOutcome{\mathrm{NNB}})\\
&=\GTenHLProbability{0.144}+\GTenHLProbability{0.144}+\GTenHLProbability{0.144}=\GTenHLProbability{0.432}.
\end{aligned}
\]
The \hyperref[subsec:t2-c1-assign-probability-to-each-value]{worked school-sale problem} completes all three PMFs.

The braces list the values in a support. A discrete support may be finite, as for the buyer count, or countably infinite, as in $\{1,2,3,\ldots\}$. A continuous measurement, such as the exact time to complete a purchase in seconds, can take values throughout an interval. Recording that time in whole seconds gives discrete values.

The \hyperref[subsec:t2-c1-notation-reference]{notation reference} shows how to write and check the sums used to build a PMF.

\secInicio[subsec:t2-c1-independent-outcomes][sec:t2-c1-introduction]

\Needspace{12\baselineskip}
\section{Notation Reference and practice procedure}\label{sec:t2-c1-notation-and-procedure}

\SectionProblemsTableOfContents{
	\SectionProblemLink{subsec:t2-c1-notation-reference}{Notation reference} \\
	\SectionProblemLink{subsec:t2-c1-practice-procedure}{Practice procedure}
}

\subsection{Notation reference}
\label{subsec:t2-c1-notation-reference}
\begin{center}
\begin{tabularx}{\linewidth}{@{}l*{3}{>{\raggedright\arraybackslash}X}@{}}
\toprule
\textbf{Symbol} & \textbf{Meaning and units} & \textbf{Example from the school shop} & \textbf{Example from the school sale} \\
\midrule
$\GTenHLOutcome{\omega}$ & One complete experimental outcome
  & \GTenHLOutcome{Debit card} is one complete outcome.
  & In $\GTenHLOutcome{\mathrm{BNN}}$, only the first visitor buys. \\
$\GTenHLOutcome{\Omega}$ & Sample space, the set of all complete outcomes
  & $\{$\GTenHLOutcome{Cash}, \GTenHLOutcome{debit card},\newline
    \phantom{$\{$}\GTenHLOutcome{credit card}, \GTenHLOutcome{mobile wallet}$\}$
  & $\{\GTenHLOutcome{\mathrm{NNN}},\GTenHLOutcome{\mathrm{BNN}},\GTenHLOutcome{\mathrm{NBN}},\GTenHLOutcome{\mathrm{NNB}},$\newline
    $\phantom{\{}\GTenHLOutcome{\mathrm{BBN}},\GTenHLOutcome{\mathrm{BNB}},\GTenHLOutcome{\mathrm{NBB}},\GTenHLOutcome{\mathrm{BBB}}\}$ \\
$\GTenHLValue{X},\GTenHLValue{x}$ & The rule that assigns a number and a particular numerical value, respectively
  & $\GTenHLValue{F}$ gives the fee in USD, and $\GTenHLValue{1}$ is one possible fee value.
  & $\GTenHLValue{X}$ counts buyers, and $\GTenHLValue{x}=\GTenHLValue{1}$ means one buyer. \\
$\GTenHLValue{X}(\GTenHLOutcome{\omega})$ & The numerical value assigned to outcome $\GTenHLOutcome{\omega}$
  & $\GTenHLValue{F}(\GTenHLOutcome{\text{debit card}})=\GTenHLValue{1}$
  & $\GTenHLValue{X}(\GTenHLOutcome{\mathrm{BNN}})=\GTenHLValue{1}$ \\
$\GTenHLValue{S}$ & Support, the set of values with positive probability
  & $\GTenHLValue{S}=\GTenHLValue{\{0,1,2\}}$ for the fee in USD.
  & $\GTenHLValue{S}=\GTenHLValue{\{0,1,2,3\}}$ for the buyer count. \\
$p_{\GTenHLValue{X}}(\GTenHLValue{x})$ & Probability assigned to value $\GTenHLValue{x}$, with no units
  & $p_{\GTenHLValue{F}}(\GTenHLValue{1})=P(\GTenHLValue{F}=\GTenHLValue{1})=\GTenHLProbability{0.5}$, the probability of a 1 USD fee.
  & $p_{\GTenHLValue{X}}(\GTenHLValue{1})=P(\GTenHLValue{X}=\GTenHLValue{1})$, the chance of exactly one buyer. \\
$\sum$ & Add the indicated terms
  & Add the probabilities of \GTenHLOutcome{debit card} and \GTenHLOutcome{credit card}.\newline
    $\sum_{\GTenHLOutcome{\omega}:\GTenHLValue{F}(\GTenHLOutcome{\omega})=\GTenHLValue{1}}P(\GTenHLOutcome{\omega})$\newline
    \({}=\GTenHLProbability{0.2}+\GTenHLProbability{0.3}=\GTenHLProbability{0.5}\)
  & $\sum_{\GTenHLValue{X}(\GTenHLOutcome{\omega})=\GTenHLValue{1}}P(\GTenHLOutcome{\omega})$ adds the probabilities of
    \GTenHLOutcome{BNN}, \GTenHLOutcome{NBN}, and \GTenHLOutcome{NNB}. \\
\bottomrule
\end{tabularx}
\end{center}

\textbf{Read an outcome and its value.} The lowercase Greek letter $\GTenHLOutcome{\omega}$ (omega) stands for one complete outcome. The capital $\GTenHLOutcome{\Omega}$ names the whole sample space. Read $\GTenHLOutcome{\omega}\in\GTenHLOutcome{\Omega}$ as \emph{omega belongs to Omega}. Read $\GTenHLValue{X}(\GTenHLOutcome{\omega})$ as \emph{X of omega}, the value the rule $\GTenHLValue{X}$ assigns to that outcome.

The \emph{probability mass function} (PMF) assigns a probability to each value of a discrete random variable. For $\GTenHLValue{X}$, it is written $p_{\GTenHLValue{X}}(\GTenHLValue{x})=P(\GTenHLValue{X}=\GTenHLValue{x})$, read ``the probability that X equals x''. The subscript identifies the random variable: $p_{\GTenHLSecond{Y}}(\GTenHLSecond{y})=P(\GTenHLSecond{Y}=\GTenHLSecond{y})$ describes revenue, while $p_{\GTenHLThird{Z}}(\GTenHLThird{z})=P(\GTenHLThird{Z}=\GTenHLThird{z})$ describes the first visitor's purchase indicator.
The expression $\GTenHLValue{X}=\GTenHLValue{x}$ describes the event that the variable takes the chosen value.

\Needspace{9\baselineskip}
\textbf{Find one PMF entry.} Fix a numerical value $\GTenHLValue{x}$ and keep every outcome that maps to it.
\[
p_{\GTenHLValue{X}}(\GTenHLValue{x})
=\sum_{\GTenHLOutcome{\omega}\in\GTenHLOutcome{\Omega}:\,\GTenHLValue{X}(\GTenHLOutcome{\omega})=\GTenHLValue{x}}P(\GTenHLOutcome{\omega}).
\]
Read the colon as \emph{such that} and the capital sigma $\sum$ as \emph{add}. The condition under the sum says which outcomes to include. Each $P(\GTenHLOutcome{\omega})$ is the probability of one complete outcome. Those outcomes are mutually exclusive, so add their probabilities once each, even when the probabilities differ. An empty sum is zero.

\textbf{Verify the PMF.} Each mass must satisfy $p_{\GTenHLValue{X}}(\GTenHLValue{x})\geq\GTenHLProbability{0}$, where $\geq$ means \emph{greater than or equal to}. The support $\GTenHLValue{S}$ contains precisely the values with positive mass. Every value outside $\GTenHLValue{S}$ has mass zero. To check \emph{normalization}, add the masses once per supported numerical value.
\[
\sum_{\GTenHLValue{x}\in\GTenHLValue{S}}p_{\GTenHLValue{X}}(\GTenHLValue{x})=1.
\]
Here $\GTenHLValue{x}\in\GTenHLValue{S}$ means \emph{x belongs to S}. The events for distinct values cannot occur together, and the supported values account for all probability. The total of 1 represents a 100\% chance of obtaining some supported value.

\textbf{Finite and countably infinite supports.} A finite sum has a last term. For a countably infinite support such as $\GTenHLValue{S}=\GTenHLValue{\{1,2,3,\ldots\}}$, the same notation represents an infinite list of non-negative terms.
\[
\sum_{\GTenHLValue{x}\in\GTenHLValue{S}}p_{\GTenHLValue{X}}(\GTenHLValue{x})=p_{\GTenHLValue{X}}(\GTenHLValue{1})+p_{\GTenHLValue{X}}(\GTenHLValue{2})+p_{\GTenHLValue{X}}(\GTenHLValue{3})+\cdots=1.
\]
Check that the partial sums, the totals of the first finitely many masses, approach 1 as more terms are included. An infinite list of possible values can have total probability 1.

See \hyperref[subsec:t2-c1-introductory-general-problem-grouping-outcomes]{the grouping example} for both sums in the same model.

\secInicio[subsec:t2-c1-notation-reference][sec:t2-c1-notation-and-procedure]

\Needspace{10\baselineskip}
\subsection{Practice procedure}
\label{subsec:t2-c1-practice-procedure}
\begin{enumerate}
\item Identify the starting point. For \textbf{given outcomes}, read their probabilities directly. For \textbf{independent events}, list the component decisions and probabilities, form the complete outcomes, and multiply within each one. For a \textbf{supplied PMF}, use its values and masses directly. Recover one missing mass by subtracting the sum of the known masses from 1.
\item Define the variable in words and give its units. If you start from outcomes, map them to values and add the probabilities of the mutually exclusive complete outcomes that give each value.
\item Check that the masses are non-negative and total 1. List the values with positive mass as the support.
\item Translate the question into an event, paying attention to endpoints. Write the formula before substituting, keep intermediate values exact, and interpret the result in context.
\end{enumerate}
\secInicio[subsec:t2-c1-practice-procedure][sec:t2-c1-notation-and-procedure]

\Needspace{13\baselineskip}
\section{General Worked Problems --- Given Outcomes or PMFs}\label{sec:t2-c1-given-outcomes-or-pmfs}

\SectionProblemsTableOfContents{
	\SectionProblemLink{subsec:t2-c1-introductory-general-problem-grouping-outcomes}{Introductory general problem: grouping outcomes} \\
	\SectionProblemLink{subsec:t2-c1-payment-methods-and-a-fee}{Payment methods and a fee} \\
	\SectionProblemLink{subsec:t2-c1-daily-insurance-claims}{Daily insurance claims} \\
	\SectionProblemLink{subsec:t2-c1-investment-results-with-a-missing-mass}{Investment results with a missing mass}
} \vspace{5mm}

When complete-outcome probabilities are given, add the masses of the mutually exclusive outcomes that map to each value. When a PMF is supplied with one missing mass, use the total of 1 to recover it.

\Needspace{8\baselineskip}
\subsection{Introductory general problem: grouping outcomes}
\label{subsec:t2-c1-introductory-general-problem-grouping-outcomes}
This problem develops the fee model from the \hyperref[subsec:t2-c1-given-outcomes]{first Introduction example}.

\subsubsection*{Mathematical wording}
An experiment has outcomes $\GTenHLOutcome{a},\GTenHLOutcome{b},\GTenHLOutcome{c},\GTenHLOutcome{d}$ with probabilities $\GTenHLProbability{0.1},\GTenHLProbability{0.2},\GTenHLProbability{0.3},\GTenHLProbability{0.4}$. Define $\GTenHLValue{X}(\GTenHLOutcome{a})=\GTenHLValue{0}$, $\GTenHLValue{X}(\GTenHLOutcome{b})=\GTenHLValue{X}(\GTenHLOutcome{c})=\GTenHLValue{1}$, $\GTenHLValue{X}(\GTenHLOutcome{d})=\GTenHLValue{2}$.

\QuestionsTasks[sec:t2-c1-given-outcomes-or-pmfs]{
	{1}/{subsubsec:t2-c1-introductory-general-problem-grouping-outcomes-solution}/
	{Construct the mapping and PMF; state the support and find $P(\GTenHLValue{X}=\GTenHLValue{1})$.}
}
\Needspace{11\baselineskip}
The table presents the same mathematical task in three financial settings.
\begin{longtable}{@{}>{\raggedright\arraybackslash}p{0.65\linewidth}>{\raggedright\arraybackslash}p{0.33\linewidth}@{}}
\toprule
\textbf{Description} & \textbf{Contextualized task} \\
\midrule
\endhead
\bottomrule
\endlastfoot
\textbf{Version 1 --- School shop processing fees.} A school shop receives one payment. There are exactly four possible payment methods: \GTenHLOutcome{cash}, \GTenHLOutcome{debit card}, \GTenHLOutcome{credit card}, or \GTenHLOutcome{a mobile wallet}. The chance of \GTenHLOutcome{cash} is \GTenHLProbability{10\%}, of \GTenHLOutcome{debit card} is \GTenHLProbability{20\%}, of \GTenHLOutcome{credit card} is \GTenHLProbability{30\%}, and of a \GTenHLOutcome{mobile wallet} is \GTenHLProbability{40\%}. The shop pays no processing fee for \GTenHLOutcome{cash}, a fee of \GTenHLValue{\$1} for either type of card, and a fee of \GTenHLValue{\$2} for a \GTenHLOutcome{mobile wallet}. These are the only processing charges.
& Match each payment method to its fee. List the possible fee amounts and their probabilities, and find the chance of a \GTenHLValue{\$1} fee. \\
\addlinespace
\textbf{Version 2 --- Bike rental late fees.} A bike rental shop charges a late fee based on when a bike is returned. A rental ends \GTenHLOutcome{on time or early} with probability \GTenHLProbability{10\%}, \GTenHLOutcome{1 to 10 minutes late} with probability \GTenHLProbability{20\%}, \GTenHLOutcome{11 to 20 minutes late} with probability \GTenHLProbability{30\%}, or \GTenHLOutcome{more than 20 minutes late} with probability \GTenHLProbability{40\%}. The shop records lateness in whole minutes. The fee is \GTenHLValue{\$0} for an on-time or early return, \GTenHLValue{\$1} for either of the two shorter delays, and \GTenHLValue{\$2} for a delay of more than 20 minutes.
& Match each return category to its fee. List the possible fee amounts and their probabilities, and find the chance of a \GTenHLValue{\$1} late fee. \\
\addlinespace
\textbf{Version 3 --- Recycling buyback payments.} At a school recycling drive, each delivery contains one type of material. The probabilities are \GTenHLProbability{10\%} for \GTenHLOutcome{non-recyclable material}, \GTenHLProbability{20\%} for \GTenHLOutcome{paper}, \GTenHLProbability{30\%} for \GTenHLOutcome{plastic}, and \GTenHLProbability{40\%} for \GTenHLOutcome{aluminium}. These are the only delivery types. A recycling centre pays the school \GTenHLValue{\$0} for a non-recyclable delivery, \GTenHLValue{\$1} for a delivery of either paper or plastic, and \GTenHLValue{\$2} for an aluminium delivery.
& Match each material to its payment amount. List the possible amounts and their probabilities, and find the chance the school receives \GTenHLValue{\$1}. \\
\end{longtable}
\subsubsection*{Complete worked solution}
\label{subsubsec:t2-c1-introductory-general-problem-grouping-outcomes-solution}
\begin{center}
\begin{tabular}{llp{0.48\linewidth}}
\toprule
Symbol & Value & Meaning in this problem \\ \midrule
$\GTenHLValue{X}$ & $\GTenHLValue{0,1,2}$ & Possible numerical values; fees in \GTenHLValue{USD} in the story \\
$p_{\GTenHLValue{X}}(\GTenHLValue{1})$ & $\GTenHLProbability{0.2}+\GTenHLProbability{0.3}$ & Mass grouped from two outcomes \\
\bottomrule
\end{tabular}
\end{center}

\par\noindent\begin{minipage}{\linewidth}
\textbf{Step 1: Define the variable.} $\GTenHLValue{X}$ is the numerical score assigned to an outcome; its units are \GTenHLValue{score units} in the abstract version. In the financial story, define $\GTenHLValue{X}$ as the \GTenHLValue{processing fee in USD}. Label \GTenHLOutcome{cash}, \GTenHLOutcome{debit card}, \GTenHLOutcome{credit card}, and \GTenHLOutcome{mobile wallet} by $\GTenHLOutcome{a},\GTenHLOutcome{b},\GTenHLOutcome{c},\GTenHLOutcome{d}$, respectively; their probabilities are $\GTenHLProbability{0.1},\GTenHLProbability{0.2},\GTenHLProbability{0.3},\GTenHLProbability{0.4}$. The fees give exactly the same numerical mapping.
\end{minipage}\par\medskip
\par\noindent\begin{minipage}{\linewidth}
\textbf{Step 2: Map outcomes.} \begin{center}
\begin{tabular}{lllll}
\toprule
\GTenHLOutcome{Outcome} & $\GTenHLOutcome{a}$ & $\GTenHLOutcome{b}$ & $\GTenHLOutcome{c}$ & $\GTenHLOutcome{d}$ \\ \midrule
\GTenHLProbability{Probability} & \GTenHLProbability{0.1} & \GTenHLProbability{0.2} & \GTenHLProbability{0.3} & \GTenHLProbability{0.4} \\
$\GTenHLValue{X}$ value & \GTenHLValue{0} & \GTenHLValue{1} & \GTenHLValue{1} & \GTenHLValue{2} \\
\bottomrule
\end{tabular}
\end{center}
The support is $\GTenHLValue{\{0,1,2\}}$. The outcomes are mutually exclusive and exhaustive; equal likelihood is not assumed.
\end{minipage}\par\medskip
\textbf{Step 3: Translate and calculate.} Apply the \hyperref[subsec:t2-c1-notation-reference]{rule for grouping outcomes}.
\begin{enumerate}
\item \textbf{Choose the value.} Fix $\GTenHLValue{x}=\GTenHLValue{1}$. In the school shop, this is a processing fee of \GTenHLValue{\$1}.
\item \textbf{Read the mapping.} For example, $\GTenHLValue{X}(\GTenHLOutcome{b})=\GTenHLValue{1}$.
\item \textbf{Keep the matching outcomes.} Keep $\GTenHLOutcome{b}$ and $\GTenHLOutcome{c}$ because both give \GTenHLValue{1}. Leave out $\GTenHLOutcome{a}$ and $\GTenHLOutcome{d}$ because they give \GTenHLValue{0} and \GTenHLValue{2}.
\item \textbf{Read their supplied probabilities.} Each entry is already the probability of a complete outcome. The mapping table gives $P(\GTenHLOutcome{b})=\GTenHLProbability{0.2}$ and $P(\GTenHLOutcome{c})=\GTenHLProbability{0.3}$.
\item \textbf{Add those probabilities.} The matching outcomes are mutually exclusive, so count each one once.
\[
p_{\GTenHLValue{X}}(\GTenHLValue{1})=P(\GTenHLValue{X}=\GTenHLValue{1})=P(\GTenHLOutcome{b})+P(\GTenHLOutcome{c})=\GTenHLProbability{0.2}+\GTenHLProbability{0.3}=\GTenHLProbability{0.5000}.
\]
The left side names the probability assigned to the value \GTenHLValue{1}; the right side shows how it is built from the outcomes. Financially, add the chances of \GTenHLOutcome{debit} and \GTenHLOutcome{credit card} payments because both produce the same \GTenHLValue{\$1} fee.
\end{enumerate}
Repeat for the other values: $p_{\GTenHLValue{X}}(\GTenHLValue{0})=P(\GTenHLOutcome{a})=\GTenHLProbability{0.1}$ and $p_{\GTenHLValue{X}}(\GTenHLValue{2})=P(\GTenHLOutcome{d})=\GTenHLProbability{0.4}$. If no outcome gives a chosen value, the sum has no terms and its value is zero; for example, $p_{\GTenHLValue{X}}(\GTenHLValue{3})=\GTenHLProbability{0}$. The complete PMF on the support is:
\begin{center}
\begin{tabular}{llll}
\toprule
$\GTenHLValue{x}$ & \GTenHLValue{0} & \GTenHLValue{1} & \GTenHLValue{2} \\ \midrule
$p_{\GTenHLValue{X}}(\GTenHLValue{x})$ & \GTenHLProbability{0.1} & \GTenHLProbability{0.5} & \GTenHLProbability{0.4} \\
\bottomrule
\end{tabular}
\end{center}

\par\medskip
\par\noindent\begin{minipage}{\linewidth}
\textbf{Step 4: Normalize and interpret.} The three masses in the PMF table are non-negative. Sum them once per value in $\GTenHLValue{S}=\GTenHLValue{\{0,1,2\}}$.
\[
\begin{aligned}
\sum_{\GTenHLValue{x}\in\GTenHLValue{S}}p_{\GTenHLValue{X}}(\GTenHLValue{x})
&=p_{\GTenHLValue{X}}(\GTenHLValue{0})+p_{\GTenHLValue{X}}(\GTenHLValue{1})+p_{\GTenHLValue{X}}(\GTenHLValue{2})\\
&=\GTenHLProbability{0.1}+\GTenHLProbability{0.5}+\GTenHLProbability{0.4}=1.
\end{aligned}
\]
The \GTenHLProbability{0.5} at $\GTenHLValue{x}=\GTenHLValue{1}$ enters this sum once. It already combines the two outcome probabilities. The score equals \GTenHLValue{1} in half of repetitions. In the school shop, the probability of a \GTenHLValue{\$1} fee is \GTenHLProbability{50\%}. The fee must be \GTenHLValue{\$0}, \GTenHLValue{\$1}, or \GTenHLValue{\$2}, so the total probability is 1. Fee amounts are measured in \GTenHLValue{USD} and the probabilities are unitless.
\end{minipage}\par\medskip
\noindent\textbf{Across the three narratives.} In Version 1, $\GTenHLValue{X}$ is the processing fee the shop pays. In Version 2, it is the late fee the renter pays. In Version 3, it is the payment the school receives for one delivery. All amounts are in \GTenHLValue{USD}. Each story gives a \GTenHLProbability{50\%} chance of a fee or payment of exactly \GTenHLValue{\$1}.\par\medskip
% Graphs follow completed calculations; captions stay with their panels.
\Needspace{6.8cm}
\par\noindent \begin{minipage}[t]{.485\linewidth}
\centering
\GTenOutcomeProbabilityGraph[width=\linewidth,height=4.7cm,xlabel={\GTenHLOutcome{Outcome}},ylabel={\GTenHLProbability{Probability}},xhighlight=GTenHLOutcome,yhighlight=GTenHLProbability,color=GTenHLProbability!32]{$a$,$b$,$c$,$d$}{.1,.2,.3,.4}
\par\small \textbf{Outcome probabilities.} Distinct outcomes have their own chances.
\end{minipage}\hfill\begin{minipage}[t]{.485\linewidth}
\centering
\GTenOutcomeValueGraph[width=\linewidth,height=4.7cm,ylabel={\GTenHLValue{Assigned score / fee (USD)}},xlabel={\GTenHLOutcome{Outcome}},xhighlight=GTenHLOutcome,yhighlight=GTenHLValue,color=GTenHLValue!32]{$a$,$b$,$c$,$d$}{0,1,1,2}
\par\small \textbf{Assigned values (Version 1).} Both $\GTenHLOutcome{b}$ and $\GTenHLOutcome{c}$ give \GTenHLValue{1}; in the financial wording these are fees in \GTenHLValue{USD}.
\end{minipage}\par\medskip
\Needspace{6.8cm}
\par\noindent \begin{minipage}[t]{.485\linewidth}
\centering
\GTenPMFGraph[width=\linewidth,height=4.7cm,xlabel={\GTenHLValue{Score / fee (USD)}},ylabel={\GTenHLProbability{Probability}},xhighlight=GTenHLValue,yhighlight=GTenHLProbability,color=GTenHLProbability!32]{0,1,2}{.1,.5,.4}
\par\small \textbf{Grouped PMF (Version 1 fee scale).} The probability at \GTenHLValue{1} is $\GTenHLProbability{0.2}+\GTenHLProbability{0.3}=\GTenHLProbability{0.5}$.
\end{minipage}\par\medskip

\secInicio[subsec:t2-c1-introductory-general-problem-grouping-outcomes][sec:t2-c1-given-outcomes-or-pmfs]

\Needspace{8\baselineskip}
\subsection{Payment methods and a fee}
\label{subsec:t2-c1-payment-methods-and-a-fee}
Use the given probabilities of complete payment outcomes with a different fee schedule.

A shop receives one payment: \GTenHLOutcome{cash} with probability \GTenHLProbability{0.25}, \GTenHLOutcome{debit} with \GTenHLProbability{0.45}, or \GTenHLOutcome{credit} with \GTenHLProbability{0.30}. \GTenHLOutcome{Cash} costs \GTenHLValue{no processing fee}; either card costs \GTenHLValue{1 USD}.

\QuestionsTasks[sec:t2-c1-given-outcomes-or-pmfs]{
	{1}/{subsubsec:t2-c1-payment-methods-and-a-fee-solution}/
	{Define the fee variable and construct its PMF. Explain why three outcomes need not give three values.}
}
\subsubsection*{Complete worked solution}
\label{subsubsec:t2-c1-payment-methods-and-a-fee-solution}

\par\noindent\begin{minipage}{\linewidth}
\textbf{Step 1: Define.} $\GTenHLValue{X}$ is the \GTenHLValue{processing fee in USD}, with support $\GTenHLValue{\{0,1\}}$. The method categories are mutually exclusive and exhaustive.
\end{minipage}\par\medskip
\par\noindent\begin{minipage}{\linewidth}
\textbf{Step 2: Map.} $\GTenHLValue{X}(\GTenHLOutcome{\text{cash}})=\GTenHLValue{0}$ and $\GTenHLValue{X}(\GTenHLOutcome{\text{debit}})=\GTenHLValue{X}(\GTenHLOutcome{\text{credit}})=\GTenHLValue{1}$. Add the probabilities of the two card outcomes.
\end{minipage}\par\medskip
\par\noindent\begin{minipage}{\linewidth}
\textbf{Step 3: Calculate.} $p_{\GTenHLValue{X}}(\GTenHLValue{0})=\GTenHLProbability{0.25}$ and $p_{\GTenHLValue{X}}(\GTenHLValue{1})=P(\GTenHLOutcome{\text{debit}})+P(\GTenHLOutcome{\text{credit}})=\GTenHLProbability{0.45}+\GTenHLProbability{0.30}=\GTenHLProbability{0.75}$. Thus $P(\GTenHLValue{X}>\GTenHLValue{0})=\GTenHLProbability{0.7500}$.
\end{minipage}\par\medskip
\par\noindent\begin{minipage}{\linewidth}
\textbf{Step 4: Check and interpret.} $\GTenHLProbability{0.25}+\GTenHLProbability{0.75}=1$. There are three method outcomes but only two fee values; a random variable need not assign different values to different outcomes.
\end{minipage}\par\medskip
% Graphs follow completed calculations; captions stay with their panels.
\Needspace{6.8cm}
\par\noindent \begin{minipage}[t]{.485\linewidth}
\centering
\GTenOutcomeProbabilityGraph[width=\linewidth,height=4.7cm,xlabel={\GTenHLOutcome{Outcome}},ylabel={\GTenHLProbability{Probability}},xhighlight=GTenHLOutcome,yhighlight=GTenHLProbability,color=GTenHLProbability!32]{cash,debit,credit}{.25,.45,.30}
\par\small \textbf{Payment-method probabilities.} The bars refer to method outcomes.
\end{minipage}\hfill\begin{minipage}[t]{.485\linewidth}
\centering
\GTenOutcomeValueGraph[width=\linewidth,height=4.7cm,ylabel={$\GTenHLValue{X}(\GTenHLOutcome{\omega})$ \GTenHLValue{(USD)}},xlabel={\GTenHLOutcome{Outcome}},xhighlight=GTenHLOutcome,yhighlight=GTenHLValue,color=GTenHLValue!32]{cash,debit,credit}{0,1,1}
\par\small \textbf{Assigned fees.} Both cards map to \GTenHLValue{1 USD}; \GTenHLOutcome{cash} maps to \GTenHLValue{zero}.
\end{minipage}\par\medskip
\Needspace{6.8cm}
\par\noindent \begin{minipage}[t]{.485\linewidth}
\centering
\GTenPMFGraph[width=\linewidth,height=4.7cm,xlabel={\GTenHLValue{$x$ (USD)}},ylabel={\GTenHLProbability{Probability}},xhighlight=GTenHLValue,yhighlight=GTenHLProbability,color=GTenHLProbability!32]{0,1}{.25,.75}
\par\small \textbf{Grouped fee PMF.} \GTenHLOutcome{Debit} and \GTenHLOutcome{credit} contribute $\GTenHLProbability{0.45}+\GTenHLProbability{0.30}=\GTenHLProbability{0.75}$ at \GTenHLValue{1 USD}.
\end{minipage}\par\medskip

\secInicio[subsec:t2-c1-payment-methods-and-a-fee][sec:t2-c1-given-outcomes-or-pmfs]

\Needspace{8\baselineskip}
\subsection{Daily insurance claims}
\label{subsec:t2-c1-daily-insurance-claims}
The supplied PMF has one missing mass. Use normalization to recover it.

The number $\GTenHLValue{X}$ of claims on a small policy group in one day has probabilities $\GTenHLProbability{0.7},\GTenHLParameter{q},\GTenHLProbability{0.1}$ for $\GTenHLValue{0,1,2}$ \GTenHLValue{claims}; no other values are possible.

\QuestionsTasks[sec:t2-c1-given-outcomes-or-pmfs]{
	{1}/{subsubsec:t2-c1-daily-insurance-claims-solution}/
	{Find $q$, verify the PMF and calculate the probability of any claim.}
}
\subsubsection*{Complete worked solution}
\label{subsubsec:t2-c1-daily-insurance-claims-solution}

\par\noindent\begin{minipage}{\linewidth}
\textbf{Step 1: Define.} $\GTenHLValue{X}$ counts \GTenHLValue{daily claims}, with proposed support $\GTenHLValue{\{0,1,2\}}$. The supplied PMF gives masses for these claim counts directly.
\end{minipage}\par\medskip
\par\noindent\begin{minipage}{\linewidth}
\textbf{Step 2: Normalize.} $\GTenHLParameter{q}=1-p_{\GTenHLValue{X}}(\GTenHLValue{0})-p_{\GTenHLValue{X}}(\GTenHLValue{2})=1-\GTenHLProbability{0.7}-\GTenHLProbability{0.1}=\GTenHLParameter{0.2}$.
\end{minipage}\par\medskip
\par\noindent\begin{minipage}{\linewidth}
\textbf{Step 3: Translate and calculate.} $P(\GTenHLValue{X}\geq\GTenHLValue{1})=p_{\GTenHLValue{X}}(\GTenHLValue{1})+p_{\GTenHLValue{X}}(\GTenHLValue{2})=\GTenHLProbability{0.2}+\GTenHLProbability{0.1}=\GTenHLProbability{0.3000}$.
\end{minipage}\par\medskip
\par\noindent\begin{minipage}{\linewidth}
\textbf{Step 4: Check and interpret.} All masses are positive and total 1. Any claim occurs with probability \GTenHLProbability{30\%}; the count itself is measured in claims, not USD.
\end{minipage}\par\medskip
% Graphs follow completed calculations; captions stay with their panels.
\noindent Only a distribution of numerical values is supplied, so the PMF below does not assume an underlying outcome list.\par\medskip
\Needspace{6.8cm}
\par\noindent \begin{minipage}[t]{.485\linewidth}
\centering
\GTenPMFGraph[width=\linewidth,height=4.7cm,xlabel={\GTenHLValue{$x$ (claims)}},ylabel={\GTenHLProbability{Probability}},xhighlight=GTenHLValue,yhighlight=GTenHLProbability,color=GTenHLProbability!32]{0,1,2}{.7,.2,.1}
\par\small \textbf{Completed claim-count PMF.} The solved parameter is $\GTenHLParameter{q}=\GTenHLParameter{0.2}$; its known mass is \GTenHLProbability{0.2}.
\end{minipage}\par\medskip

\secInicio[subsec:t2-c1-daily-insurance-claims][sec:t2-c1-given-outcomes-or-pmfs]

\Needspace{8\baselineskip}
\subsection{Investment results with a missing mass}
\label{subsec:t2-c1-investment-results-with-a-missing-mass}
The supplied PMF has one missing mass. Use normalization to test the proposed value and recover the correct mass.

A hypothetical one-month net gain $\GTenHLValue{G}$ takes values $\GTenHLValue{-20,0,30}$ \GTenHLValue{USD} with masses $\GTenHLProbability{0.2},\GTenHLProbability{0.5},\GTenHLParameter{q}$. A student proposes $\GTenHLParameter{q}=\GTenHLParameter{0.4}$.

\QuestionsTasks[sec:t2-c1-given-outcomes-or-pmfs]{
	{1}/{subsubsec:t2-c1-investment-results-with-a-missing-mass-solution}/
	{Correct the PMF, identify the support and calculate the probability of a non-negative gain.}
}
\subsubsection*{Complete worked solution}
\label{subsubsec:t2-c1-investment-results-with-a-missing-mass-solution}

\par\noindent\begin{minipage}{\linewidth}
\textbf{Step 1: Define.} $\GTenHLValue{G}$ is monetary net gain, not a count. A finite set of monetary outcomes is discrete even though money is involved.
\end{minipage}\par\medskip
\par\noindent\begin{minipage}{\linewidth}
\textbf{Step 2: Normalize.} $\GTenHLParameter{q}=1-\GTenHLProbability{0.2}-\GTenHLProbability{0.5}=\GTenHLParameter{0.3}$; the proposed \GTenHLParameter{0.4} would make total mass \GTenHLProbability{1.1}, so it is invalid. Support is $\GTenHLValue{\{-20,0,30\}}$ \GTenHLValue{USD}.
\end{minipage}\par\medskip
\par\noindent\begin{minipage}{\linewidth}
\textbf{Step 3: Calculate.} $P(\GTenHLValue{G}\geq\GTenHLValue{0})=p_{\GTenHLValue{G}}(\GTenHLValue{0})+p_{\GTenHLValue{G}}(\GTenHLValue{30})=\GTenHLProbability{0.5}+\GTenHLProbability{0.3}=\GTenHLProbability{0.8000}$.
\end{minipage}\par\medskip
\par\noindent\begin{minipage}{\linewidth}
\textbf{Step 4: Check and interpret.} The corrected masses are non-negative and total 1. A loss remains possible with probability \GTenHLProbability{0.2}. Zero gain is included in non-negative gain.
\end{minipage}\par\medskip
% Graphs follow completed calculations; captions stay with their panels.
\noindent Only a distribution of numerical values is supplied, so the PMF below does not assume an underlying outcome list.\par\medskip
\Needspace{6.8cm}
\par\noindent \begin{minipage}[t]{.485\linewidth}
\centering
\GTenPMFGraph[width=\linewidth,height=4.7cm,xlabel={\GTenHLValue{$g$ (USD)}},xstep=10,ylabel={\GTenHLProbability{Probability}},xhighlight=GTenHLValue,yhighlight=GTenHLProbability,color=GTenHLProbability!32]{-20,0,30}{.2,.5,.3}
\par\small \textbf{Corrected gain PMF.} The solved parameter is $\GTenHLParameter{q}=\GTenHLParameter{0.3}$; the bar at \GTenHLValue{30 USD} uses known mass \GTenHLProbability{0.3}. Gain values retain their numerical spacing.
\end{minipage}\par\medskip

\secInicio[subsec:t2-c1-investment-results-with-a-missing-mass][sec:t2-c1-given-outcomes-or-pmfs]

\Needspace{13\baselineskip}
\section{General Worked Problems --- Independent Events}\label{sec:t2-c1-independent-events}

\SectionProblemsTableOfContents{
	\SectionProblemLink{subsec:t2-c1-assign-probability-to-each-value}{School sale: three probability mass functions} \\
	\SectionProblemLink{subsec:t2-c1-notebook-buyers}{Notebook buyers}
} \vspace{5mm}

Multiply the probabilities of independent decisions within each complete outcome. Then add the probabilities of the mutually exclusive complete outcomes that map to each value.

\Needspace{8\baselineskip}
\subsection{School sale: three probability mass functions}
\label{subsec:t2-c1-assign-probability-to-each-value}

Return to the \hyperref[subsec:t2-c1-independent-outcomes]{second Introduction example}, where we built the ordered outcomes for the three-visitor school sale. Use the buyer count $\GTenHLValue{X}$, the revenue $\GTenHLSecond{Y}$, and the first-visitor indicator $\GTenHLThird{Z}$ from the Introduction.

For each visitor in the school sale, assign
\[
P(\GTenHLOutcome{\mathrm{B}})=\GTenHLProbability{0.4},\qquad P(\GTenHLOutcome{\mathrm{N}})=\GTenHLProbability{0.6},
\]
where B means buying a notebook and N means not buying. These probabilities add to 1. Assume that the three visitors decide independently and that each has these same probabilities. Independence means that one visitor's decision does not change the probabilities for another visitor.

\QuestionsTasks[sec:t2-c1-independent-events]{
	{1}/{subsubsec:t2-c1-school-sale-solution}/
	{Find the probabilities of the eight complete outcomes. Use them to build and check all three PMFs.}
}

\subsubsection*{Complete worked solution}
\label{subsubsec:t2-c1-school-sale-solution}
\textbf{Step 1: Find the probability of each complete outcome.} As in the Introduction, independence allows us to multiply the probabilities of the three decisions within each ordered outcome. For example,
\[
\begin{aligned}
P(\GTenHLOutcome{\mathrm{BNN}})&=\GTenHLProbability{0.4}\times\GTenHLProbability{0.6}\times\GTenHLProbability{0.6}=\GTenHLProbability{0.144},\\
P(\GTenHLOutcome{\mathrm{BBN}})&=\GTenHLProbability{0.4}\times\GTenHLProbability{0.4}\times\GTenHLProbability{0.6}=\GTenHLProbability{0.096}.
\end{aligned}
\]
Applying the same rule to all eight outcomes gives
\[
\begin{aligned}
	P(\GTenHLOutcome{\mathrm{NNN}})&=\GTenHLProbability{0.6}^3=\GTenHLProbability{0.216}, & P(\GTenHLOutcome{\mathrm{BNN}})&=\GTenHLProbability{0.4}\times\GTenHLProbability{0.6}^2=\GTenHLProbability{0.144},\\
	P(\GTenHLOutcome{\mathrm{NBN}})&=\GTenHLProbability{0.6}\times\GTenHLProbability{0.4}\times\GTenHLProbability{0.6}=\GTenHLProbability{0.144}, & P(\GTenHLOutcome{\mathrm{NNB}})&=\GTenHLProbability{0.6}^2\times\GTenHLProbability{0.4}=\GTenHLProbability{0.144},\\
	P(\GTenHLOutcome{\mathrm{BBN}})&=\GTenHLProbability{0.4}^2\times\GTenHLProbability{0.6}=\GTenHLProbability{0.096}, & P(\GTenHLOutcome{\mathrm{BNB}})&=\GTenHLProbability{0.4}\times\GTenHLProbability{0.6}\times\GTenHLProbability{0.4}=\GTenHLProbability{0.096},\\
	P(\GTenHLOutcome{\mathrm{NBB}})&=\GTenHLProbability{0.6}\times\GTenHLProbability{0.4}^2=\GTenHLProbability{0.096}, & P(\GTenHLOutcome{\mathrm{BBB}})&=\GTenHLProbability{0.4}^3=\GTenHLProbability{0.064}.
\end{aligned}
\]
Several complete outcomes share a probability, and the probabilities vary with the number of buyers. Check the total by adding the probabilities of the eight mutually exclusive outcomes.
\[
\GTenHLProbability{0.216}+3(\GTenHLProbability{0.144})+3(\GTenHLProbability{0.096})+\GTenHLProbability{0.064}=\GTenHLProbability{1}.
\]

\textbf{Step 2: Group outcomes by the value of each variable.} For each value, add the probabilities of the mutually exclusive complete outcomes that give it. Each product in Step 1 describes one sequence of independent decisions. Each sum here combines alternative sequences.

\begin{itemize}
	\item \textbf{Example 1 (Number of buyers).}
	The outcomes \GTenHLOutcome{BNN}, \GTenHLOutcome{NBN}, and \GTenHLOutcome{NNB} produce $\GTenHLValue{X}=\GTenHLValue{1}$, so
	\[
	\begin{aligned}
		p_{\GTenHLValue{X}}(\GTenHLValue{1})
		&=P(\GTenHLOutcome{\mathrm{BNN}})+P(\GTenHLOutcome{\mathrm{NBN}})+P(\GTenHLOutcome{\mathrm{NNB}})\\
		&=\GTenHLProbability{0.144}+\GTenHLProbability{0.144}+\GTenHLProbability{0.144}=\GTenHLProbability{0.432}.
	\end{aligned}
	\]
	Similarly, \GTenHLOutcome{BBN}, \GTenHLOutcome{BNB}, and \GTenHLOutcome{NBB} produce $\GTenHLValue{X}=\GTenHLValue{2}$, so
	\[
	p_{\GTenHLValue{X}}(\GTenHLValue{2})=\GTenHLProbability{0.096}+\GTenHLProbability{0.096}+\GTenHLProbability{0.096}=\GTenHLProbability{0.288}.
	\]
	Only \GTenHLOutcome{NNN} gives $\GTenHLValue{X}=\GTenHLValue{0}$, and only \GTenHLOutcome{BBB} gives $\GTenHLValue{X}=\GTenHLValue{3}$.
	The complete PMF is
	\[
	\begin{aligned}
		p_{\GTenHLValue{X}}(\GTenHLValue{0})&=\GTenHLProbability{0.216}, & p_{\GTenHLValue{X}}(\GTenHLValue{1})&=\GTenHLProbability{0.432},\\
		p_{\GTenHLValue{X}}(\GTenHLValue{2})&=\GTenHLProbability{0.288}, & p_{\GTenHLValue{X}}(\GTenHLValue{3})&=\GTenHLProbability{0.064}.
	\end{aligned}
	\]

	\item \textbf{Example 2 (Sales revenue).}
	Each notebook sells for 5 USD, so $\GTenHLSecond{Y}=5\GTenHLValue{X}$ when revenue
	is recorded in USD. The same outcomes that produce one
	buyer produce revenue of 5 USD. Thus the complete PMF is
	\[
	\begin{aligned}
		p_{\GTenHLSecond{Y}}(\GTenHLSecond{0})&=\GTenHLProbability{0.216}, & p_{\GTenHLSecond{Y}}(\GTenHLSecond{5})&=\GTenHLProbability{0.432},\\
		p_{\GTenHLSecond{Y}}(\GTenHLSecond{10})&=\GTenHLProbability{0.288}, & p_{\GTenHLSecond{Y}}(\GTenHLSecond{15})&=\GTenHLProbability{0.064}.
	\end{aligned}
	\]
	The revenue values have monetary units; their
	probabilities do not.

	\item \textbf{Example 3 (First visitor's purchase).}
	The first visitor buys in \GTenHLOutcome{BNN}, \GTenHLOutcome{BBN}, \GTenHLOutcome{BNB}, and \GTenHLOutcome{BBB}, so
	\[
	\begin{aligned}
		p_{\GTenHLThird{Z}}(\GTenHLThird{1})
		&=\GTenHLProbability{0.144}+\GTenHLProbability{0.096}+\GTenHLProbability{0.096}+\GTenHLProbability{0.064}\\
		&=\GTenHLProbability{0.400}.
	\end{aligned}
	\]
	The first visitor does not buy in \GTenHLOutcome{NNN}, \GTenHLOutcome{NBN}, \GTenHLOutcome{NNB},
	and \GTenHLOutcome{NBB}, so
	\[
	\begin{aligned}
		p_{\GTenHLThird{Z}}(\GTenHLThird{0})
		&=\GTenHLProbability{0.216}+\GTenHLProbability{0.144}+\GTenHLProbability{0.144}+\GTenHLProbability{0.096}\\
		&=\GTenHLProbability{0.600}.
	\end{aligned}
	\]
\end{itemize}

\textbf{Step 3: Check and interpret.} For $\GTenHLValue{X}$ and $\GTenHLSecond{Y}$, the probabilities total $\GTenHLProbability{0.216}+\GTenHLProbability{0.432}+\GTenHLProbability{0.288}+\GTenHLProbability{0.064}=1$; for $\GTenHLThird{Z}$, they total $\GTenHLProbability{0.600}+\GTenHLProbability{0.400}=1$. All entries are non-negative, and values outside each support have mass zero. The probabilities for $\GTenHLThird{Z}$ agree with the original chances of the first visitor not buying or buying. Although $\GTenHLValue{X}$ and $\GTenHLSecond{Y}$ have different supports and units, corresponding values have the same probabilities because $\GTenHLSecond{Y}=5\GTenHLValue{X}$. The variable $\GTenHLThird{Z}$ groups the outcomes differently.

% Graphs follow completed calculations; captions stay with their panels.
\Needspace{6.8cm}
\par\noindent \begin{minipage}[t]{.485\linewidth}
\centering
\GTenOutcomeProbabilityGraph[width=\linewidth,height=4.7cm,rotate=35,ystep=.05,xlabel={\GTenHLOutcome{Outcome}},ylabel={\GTenHLProbability{Probability}},xhighlight=GTenHLOutcome,yhighlight=GTenHLProbability,color=GTenHLProbability!32]{NNN,BNN,NBN,NNB,BBN,BNB,NBB,BBB}{.216,.144,.144,.144,.096,.096,.096,.064}
\par\small \textbf{School sale: outcome probabilities.} Each bar describes one complete three-visitor outcome.
\end{minipage}\hfill\begin{minipage}[t]{.485\linewidth}
\centering
\GTenOutcomeValueGraph[width=\linewidth,height=4.7cm,rotate=35,ylabel={$\GTenHLValue{X}(\GTenHLOutcome{\omega})$ \GTenHLValue{(buyers)}},xlabel={\GTenHLOutcome{Outcome}},xhighlight=GTenHLOutcome,yhighlight=GTenHLValue,color=GTenHLValue!32]{NNN,BNN,NBN,NNB,BBN,BNB,NBB,BBB}{0,1,1,1,2,2,2,3}
\par\small \textbf{Assigned buyer counts.} The height is $\GTenHLValue{X}(\GTenHLOutcome{\omega})$ in \GTenHLValue{buyers}, not a probability.
\end{minipage}\par\medskip
\Needspace{6.8cm}
\par\noindent \begin{minipage}[t]{.485\linewidth}
\centering
\GTenOutcomeValueGraph[width=\linewidth,height=4.7cm,rotate=35,ystep=5,ylabel={$\GTenHLSecond{Y}(\GTenHLOutcome{\omega})$ \GTenHLSecond{(USD)}},xlabel={\GTenHLOutcome{Outcome}},xhighlight=GTenHLOutcome,yhighlight=GTenHLSecond,color=GTenHLSecond!32]{NNN,BNN,NBN,NNB,BBN,BNB,NBB,BBB}{0,5,5,5,10,10,10,15}
\par\small \textbf{Assigned revenue.} The height is $\GTenHLSecond{Y}(\GTenHLOutcome{\omega})=5\GTenHLValue{X}(\GTenHLOutcome{\omega})$ in \GTenHLSecond{USD}.
\end{minipage}\hfill\begin{minipage}[t]{.485\linewidth}
\centering
\GTenOutcomeValueGraph[width=\linewidth,height=4.7cm,rotate=35,ylabel={$\GTenHLThird{Z}(\GTenHLOutcome{\omega})$ \GTenHLThird{(indicator)}},xlabel={\GTenHLOutcome{Outcome}},xhighlight=GTenHLOutcome,yhighlight=GTenHLThird,color=GTenHLThird!32]{NNN,BNN,NBN,NNB,BBN,BNB,NBB,BBB}{0,1,0,0,1,1,0,1}
\par\small \textbf{Assigned indicator.} $\GTenHLThird{Z}(\GTenHLOutcome{\omega})$ is \GTenHLThird{1} exactly when the first visitor buys; it has no units.
\end{minipage}\par\medskip
\Needspace{6.8cm}
\par\noindent \begin{minipage}[t]{.485\linewidth}
\centering
\GTenPMFGraph[width=\linewidth,height=4.7cm,xlabel={\GTenHLValue{$x$ (buyers)}},ylabel={\GTenHLProbability{Probability}},xhighlight=GTenHLValue,yhighlight=GTenHLProbability,color=GTenHLProbability!32]{0,1,2,3}{.216,.432,.288,.064}
\par\small \textbf{Buyer-count PMF.} Group outcomes by $\GTenHLValue{X}$ and add their probabilities.
\end{minipage}\hfill\begin{minipage}[t]{.485\linewidth}
\centering
\GTenPMFGraph[width=\linewidth,height=4.7cm,xlabel={\GTenHLSecond{$y$ (USD)}},xstep=5,ylabel={\GTenHLProbability{Probability}},xhighlight=GTenHLSecond,yhighlight=GTenHLProbability,color=GTenHLProbability!32]{0,5,10,15}{.216,.432,.288,.064}
\par\small \textbf{Revenue PMF.} Group outcomes by $\GTenHLSecond{Y}$; probabilities have no USD units.
\end{minipage}\par\medskip
\Needspace{6.8cm}
\par\noindent \begin{minipage}[t]{.485\linewidth}
\centering
\GTenPMFGraph[width=\linewidth,height=4.7cm,xlabel={\GTenHLThird{$z$ (indicator)}},ylabel={\GTenHLProbability{Probability}},xhighlight=GTenHLThird,yhighlight=GTenHLProbability,color=GTenHLProbability!32]{0,1}{.6,.4}
\par\small \textbf{First-visitor PMF.} Group outcomes by $\GTenHLThird{Z}$; these bars give probabilities of indicator values.
\end{minipage}\par\medskip

\textbf{Check.} Is $p_{\GTenHLSecond{Y}}(\GTenHLSecond{5})=\GTenHLProbability{0.432}$ a payment of \GTenHLProbability{0.432} USD? What does it mean?
\par\nopagebreak
\emph{Suggested answer:} No. It is a dimensionless probability: there is a \GTenHLProbability{43.2\%} chance that revenue equals \GTenHLSecond{5 USD}.

\textbf{Check.} Are the events $\GTenHLValue{X}=\GTenHLValue{1}$, $\GTenHLSecond{Y}=\GTenHLSecond{5}$, and $\GTenHLThird{Z}=\GTenHLThird{1}$ the same?
\par\nopagebreak
\emph{Suggested answer:} The first two are the same event, consisting of \GTenHLOutcome{BNN}, \GTenHLOutcome{NBN}, and \GTenHLOutcome{NNB}. The event $\GTenHLThird{Z}=\GTenHLThird{1}$ consists of \GTenHLOutcome{BNN}, \GTenHLOutcome{BBN}, \GTenHLOutcome{BNB}, and \GTenHLOutcome{BBB}, so it is different. Their probabilities are \GTenHLProbability{0.432}, \GTenHLProbability{0.432}, and \GTenHLProbability{0.400}, respectively.

\secInicio[subsec:t2-c1-assign-probability-to-each-value][sec:t2-c1-independent-events]

\Needspace{8\baselineskip}
\subsection{Notebook buyers}
\label{subsec:t2-c1-notebook-buyers}
This time, there are two visitors and a different price per sale.

Two visitors buy independently with probability $\GTenHLProbability{1/2}$ each. Let \GTenHLOutcome{B} mean a purchase and \GTenHLOutcome{N} no purchase. Each sale earns \GTenHLSecond{4 USD}.

\QuestionsTasks[sec:t2-c1-independent-events]{
	{1}/{subsubsec:t2-c1-notebook-buyers-solution}/
	{{List outcomes, define the buyer count $X$, build its PMF and find the chance of revenue $4$ USD.}}
}
\subsubsection*{Complete worked solution}
\label{subsubsec:t2-c1-notebook-buyers-solution}

\par\noindent\begin{minipage}{\linewidth}
\textbf{Step 1: Define and list.} $\GTenHLValue{X}$ counts buyers, with support $\GTenHLValue{\{0,1,2\}}$. Revenue is $\GTenHLSecond{Y}=4\GTenHLValue{X}$ \GTenHLSecond{USD}, with support $\GTenHLSecond{\{0,4,8\}}$ \GTenHLSecond{USD}.
\end{minipage}\par\medskip
\par\noindent\begin{minipage}{\linewidth}
\textbf{Step 2: Check the model and map.} Multiply the two decision probabilities within each ordered outcome because the decisions are independent. Both decisions have probability \GTenHLProbability{1/2} for each visitor, so \GTenHLOutcome{BB}, \GTenHLOutcome{BN}, \GTenHLOutcome{NB}, and \GTenHLOutcome{NN} each have probability $(\GTenHLProbability{1/2})^2=\GTenHLProbability{1/4}$. Their $\GTenHLValue{X}$ values are \GTenHLValue{2, 1, 1, 0}.
\end{minipage}\par\medskip
\par\noindent\begin{minipage}{\linewidth}
\textbf{Step 3: Calculate.} Add the probabilities of the mutually exclusive complete outcomes with one buyer. Thus $p_{\GTenHLValue{X}}(\GTenHLValue{1})=P(\GTenHLOutcome{\mathrm{BN}})+P(\GTenHLOutcome{\mathrm{NB}})=\GTenHLProbability{1/4}+\GTenHLProbability{1/4}=\GTenHLProbability{0.5}$.\begin{center}
\begin{tabular}{llll}
\toprule
$\GTenHLValue{x}$ \GTenHLValue{(buyers)} & \GTenHLValue{0} & \GTenHLValue{1} & \GTenHLValue{2} \\ \midrule
$p_{\GTenHLValue{X}}(\GTenHLValue{x})$ & \GTenHLProbability{0.25} & \GTenHLProbability{0.50} & \GTenHLProbability{0.25} \\
\bottomrule
\end{tabular}
\end{center}
$P(\GTenHLSecond{Y}=\GTenHLSecond{4})=P(\GTenHLValue{X}=\GTenHLValue{1})=\GTenHLProbability{0.5000}$.
\end{minipage}\par\medskip
\par\noindent\begin{minipage}{\linewidth}
\textbf{Step 4: Check and interpret.} The table totals 1. A revenue of \GTenHLSecond{4 USD} occurs in two of the four equally likely outcomes; \GTenHLProbability{0.5} is a probability, not revenue.
\end{minipage}\par\medskip
% Graphs follow completed calculations; captions stay with their panels.
\Needspace{6.8cm}
\par\noindent \begin{minipage}[t]{.485\linewidth}
\centering
\GTenOutcomeProbabilityGraph[width=\linewidth,height=4.7cm,xlabel={\GTenHLOutcome{Outcome}},ylabel={\GTenHLProbability{Probability}},xhighlight=GTenHLOutcome,yhighlight=GTenHLProbability,color=GTenHLProbability!32]{BB,BN,NB,NN}{.25,.25,.25,.25}
\par\small \textbf{Complete-outcome probabilities.} Each outcome has the same product of two decision probabilities.
\end{minipage}\hfill\begin{minipage}[t]{.485\linewidth}
\centering
\GTenOutcomeValueGraph[width=\linewidth,height=4.7cm,ylabel={$\GTenHLValue{X}(\GTenHLOutcome{\omega})$ \GTenHLValue{(buyers)}},xlabel={\GTenHLOutcome{Outcome}},xhighlight=GTenHLOutcome,yhighlight=GTenHLValue,color=GTenHLValue!32]{BB,BN,NB,NN}{2,1,1,0}
\par\small \textbf{Assigned buyer counts.} The bars show $\GTenHLValue{X}(\GTenHLOutcome{\omega})$, measured in \GTenHLValue{buyers}.
\end{minipage}\par\medskip
\Needspace{6.8cm}
\par\noindent \begin{minipage}[t]{.485\linewidth}
\centering
\GTenOutcomeValueGraph[width=\linewidth,height=4.7cm,ystep=4,ylabel={$\GTenHLSecond{Y}(\GTenHLOutcome{\omega})$ \GTenHLSecond{(USD)}},xlabel={\GTenHLOutcome{Outcome}},xhighlight=GTenHLOutcome,yhighlight=GTenHLSecond,color=GTenHLSecond!32]{BB,BN,NB,NN}{8,4,4,0}
\par\small \textbf{Assigned revenue.} Each \GTenHLOutcome{outcome} maps to $\GTenHLSecond{Y}(\GTenHLOutcome{\omega})=4\GTenHLValue{X}(\GTenHLOutcome{\omega})$ in \GTenHLSecond{USD}; these heights are assigned values.
\end{minipage}\par\medskip

\Needspace{6.8cm}
\par\noindent \begin{minipage}[t]{.485\linewidth}
\centering
\GTenPMFGraph[width=\linewidth,height=4.7cm,xlabel={\GTenHLValue{$x$ (buyers)}},ylabel={\GTenHLProbability{Probability}},xhighlight=GTenHLValue,yhighlight=GTenHLProbability,color=GTenHLProbability!32]{0,1,2}{.25,.5,.25}
\par\small \textbf{Buyer-count PMF.} \GTenHLOutcome{BN} and \GTenHLOutcome{NB} are grouped at \GTenHLValue{one buyer}.
\end{minipage}\hfill\begin{minipage}[t]{.485\linewidth}
\centering
\GTenPMFGraph[width=\linewidth,height=4.7cm,xlabel={\GTenHLSecond{$y$ (USD)}},xstep=4,ylabel={\GTenHLProbability{Probability}},xhighlight=GTenHLSecond,yhighlight=GTenHLProbability,color=GTenHLProbability!32]{0,4,8}{.25,.5,.25}
\par\small \textbf{Revenue PMF.} $\GTenHLSecond{Y}=4\GTenHLValue{X}$ in \GTenHLSecond{USD}; the corresponding probabilities are unchanged.
\end{minipage}\par\medskip

\secInicio[subsec:t2-c1-notebook-buyers][sec:t2-c1-independent-events]

\Needspace{13\baselineskip}
\section{Advanced Problems}\label{sec:t2-c1-advanced}

\SectionProblemsTableOfContents{
	\SectionProblemLink{subsec:t2-c1-advanced-general-problem-a-finite-model}{Advanced general problem: a finite model} \\
	\SectionProblemLink{subsec:t2-c1-advanced-general-problem-an-infinite-model}{Advanced general problem: an infinite model} \\
	\SectionProblemLink{subsec:t2-c1-calls-until-a-sale}{Calls until a sale} \\
	\SectionProblemLink{subsec:t2-c1-delivery-records-two-numerical-descriptions}{Delivery records: two numerical descriptions}
}

\Needspace{8\baselineskip}
\subsection{Advanced general problem: a finite model}
\label{subsec:t2-c1-advanced-general-problem-a-finite-model}
\subsubsection*{Mathematical wording}
The masses in this model are given in proportion to one another. Use normalization to determine their common scale.

A model assigns masses $\GTenHLParameter{c},\GTenHLParameter{2c},\GTenHLParameter{3c}$ to scores $\GTenHLValue{0,2,5}$, respectively; let $\GTenHLValue{X}$ be its score.

\QuestionsTasks[sec:t2-c1-advanced]{
	{1}/{subsubsec:t2-c1-advanced-general-problem-a-finite-model-solution}/
	{Determine $\GTenHLParameter{c}$, verify the PMF, and calculate $P(\GTenHLValue{X}\geq\GTenHLValue{2})$.}
}
\Needspace{11\baselineskip}
The table presents the same mathematical task in three financial settings.
\begin{longtable}{@{}>{\raggedright\arraybackslash}p{0.55\linewidth}>{\raggedright\arraybackslash}p{0.41\linewidth}@{}}
\toprule
\textbf{Description} & \textbf{Contextualized task} \\
\midrule
\endhead
\bottomrule
\endlastfoot
\textbf{Version 1 --- Promotional refunds.} A shop offers a promotional refund on one purchase. The refund can only be \GTenHLValue{\$0}, \GTenHLValue{\$2}, or \GTenHLValue{\$5}. A \GTenHLValue{\$2} refund is twice as likely as no refund, and a \GTenHLValue{\$5} refund is three times as likely as no refund. These are the only possible refund amounts.
& Find the probability of no refund and use it to find the other refund probabilities. Verify that all are non-negative and sum to 1, then find the chance of a refund of at least \GTenHLValue{\$2}. \\
\addlinespace
\textbf{Version 2 --- Fundraising draw prizes.} At a school fundraiser, a participant draws one ticket from a box. A ticket pays \GTenHLValue{\$0}, \GTenHLValue{\$2}, or \GTenHLValue{\$5}. The box contains twice as many \GTenHLValue{\$2} tickets as no-prize tickets and three times as many \GTenHLValue{\$5} tickets as no-prize tickets. Each ticket has the same chance of being drawn.
& Find the probability of drawing a no-prize ticket and use it to find the other payout probabilities. Verify that all are non-negative and sum to 1, then find the chance of winning at least \GTenHLValue{\$2}. \\
\addlinespace
\textbf{Version 3 --- Equipment maintenance costs.} A student printing business sets aside money for printer maintenance. One inspection can lead to no repair at a cost of \GTenHLValue{\$0}, an adjustment costing \GTenHLValue{\$2}, or a replacement part costing \GTenHLValue{\$5}. Exactly one result occurs per inspection. An adjustment is twice as likely as no repair, and a replacement is three times as likely as no repair.
& Find the probability of no repair and use it to find the other cost probabilities. Verify that all are non-negative and sum to 1, then find the chance of a maintenance cost of at least \GTenHLValue{\$2}. \\
\end{longtable}
\subsubsection*{Complete worked solution}
\label{subsubsec:t2-c1-advanced-general-problem-a-finite-model-solution}
\par\noindent\begin{minipage}{\linewidth}
\textbf{Step 1: Define and identify support.} $\GTenHLValue{X}$ is the score with support $\GTenHLValue{\{0,2,5\}}$. In the financial story, $\GTenHLValue{X}$ is the \GTenHLValue{refund in USD}. If $\GTenHLParameter{c}$ is the chance of no refund, the relative likelihoods give masses $\GTenHLParameter{c},\GTenHLParameter{2c},\GTenHLParameter{3c}$ for refunds of \GTenHLValue{0, 2, and 5 USD}. This is the same finite model; normalization must be checked.
\end{minipage}\par\medskip
\par\noindent\begin{minipage}{\linewidth}
\textbf{Step 2: Normalize.} $\sum_{\GTenHLValue{x}}p_{\GTenHLValue{X}}(\GTenHLValue{x})=\GTenHLParameter{c}+\GTenHLParameter{2c}+\GTenHLParameter{3c}=\GTenHLParameter{6c}=1$, so $\GTenHLParameter{c}=\GTenHLParameter{1/6}$. The masses are $\GTenHLProbability{1/6},\GTenHLProbability{1/3},\GTenHLProbability{1/2}$, all positive and with total 1, so they define a PMF. In the story, these are the chances of refunds of \GTenHLValue{\$0}, \GTenHLValue{\$2}, and \GTenHLValue{\$5}, respectively.
\end{minipage}\par\medskip
\par\noindent\begin{minipage}{\linewidth}
\textbf{Step 3: Translate and calculate.} $P(\GTenHLValue{X}\geq\GTenHLValue{2})=p_{\GTenHLValue{X}}(\GTenHLValue{2})+p_{\GTenHLValue{X}}(\GTenHLValue{5})=\GTenHLProbability{1/3}+\GTenHLProbability{1/2}=\GTenHLProbability{5/6}\approx\GTenHLProbability{0.8333}$.
\end{minipage}\par\medskip
\par\noindent\begin{minipage}{\linewidth}
\textbf{Step 4: Check and interpret.} The probability lies between 0 and 1. The model has only three possible values despite the gap from \GTenHLValue{2} to \GTenHLValue{5}. Financially, the chance of a refund of at least \GTenHLValue{\$2} is approximately \GTenHLProbability{83.33\%}. The refund is measured in \GTenHLValue{USD}; the probability has no units.
\end{minipage}\par\medskip
\noindent\textbf{Across the three narratives.} The variable $\GTenHLValue{X}$ measures the refund in Version 1, the prize payout in Version 2, and the maintenance cost in Version 3, all in \GTenHLValue{USD}. In each case, the chance that the respective refund, prize, or cost is at least \GTenHLValue{\$2} is $\GTenHLProbability{5/6}$, or about \GTenHLProbability{83.33\%}.\par\medskip
% Graphs follow completed calculations; captions stay with their panels.
\noindent Only a distribution of numerical values is supplied, so the PMF below does not assume an underlying outcome list.\par\medskip
\Needspace{6.8cm}
\par\noindent \begin{minipage}[t]{.485\linewidth}
\centering
\GTenPMFGraph[width=\linewidth,height=4.7cm,xlabel={\GTenHLValue{Score / refund (USD)}},ylabel={\GTenHLProbability{Probability}},xhighlight=GTenHLValue,yhighlight=GTenHLProbability,color=GTenHLProbability!32]{0,2,5}{1/6,1/3,1/2}
\par\small \textbf{Normalized finite PMF (Version 1 refund scale).} Only \GTenHLValue{0}, \GTenHLValue{2}, and \GTenHLValue{5} are supported; the gap is preserved. Heights display rounded probabilities.
\end{minipage}\par\medskip

\secInicio[subsec:t2-c1-advanced-general-problem-a-finite-model][sec:t2-c1-advanced]

\Needspace{8\baselineskip}
\subsection{Advanced general problem: an infinite model}
\label{subsec:t2-c1-advanced-general-problem-an-infinite-model}
\subsubsection*{Mathematical wording}
A model assigns $p_{\GTenHLValue{T}}(\GTenHLValue{t})=\GTenHLProbability{2^{-t}}$ for $\GTenHLValue{t}=\GTenHLValue{1,2,\ldots}$, where $\GTenHLValue{T}$ is a \GTenHLValue{positive-integer trial count}.

\QuestionsTasks[sec:t2-c1-advanced]{
	{1}/{subsubsec:t2-c1-advanced-general-problem-an-infinite-model-solution}/
	{Verify the PMF and calculate $P(\GTenHLValue{T}>\GTenHLValue{3})$.}
}
\Needspace{28\baselineskip}
The table presents the same mathematical task in three financial settings.
\begin{longtable}{@{}>{\raggedright\arraybackslash}p{0.55\linewidth}>{\raggedright\arraybackslash}p{0.41\linewidth}@{}}
\toprule
\textbf{Description} & \textbf{Contextualized task} \\
\midrule
\endhead
\bottomrule
\endlastfoot
\textbf{Version 1 --- Calls until the first sale.} A student running a school enterprise calls potential customers until the first sale. Each call is an independent attempt, with a constant one-half chance of making a sale. Count every call, including the call that makes the first sale. There is no limit on the number of calls allowed, and calling stops at the first sale.
& Verify that the probabilities for all possible call counts are non-negative and sum to 1 over arbitrarily many calls. Find the chance that the first sale takes \GTenHLValue{more than 3 calls}. \\
\addlinespace
\textbf{Version 2 --- Finding a quote within budget.} A student business requests supplier quotes one at a time and stops at the first quote within its budget. Each quote has a \GTenHLProbability{one-half} chance of meeting the budget, independent of the other quotes. There is no fixed limit on the number of quotes. Count each quote received, including the first one within budget.
& Verify that the probabilities for all possible quote counts are non-negative and sum to 1 over arbitrarily many quotes. Find the chance that the first quote within budget takes \GTenHLValue{more than 3 quotes}. \\
\addlinespace
\textbf{Version 3 --- Finding an invoice with a discount.} A bookkeeping trainee samples past supplier invoices to find an example of an early-payment discount. Exactly half of the invoices qualify for the discount. On each draw the trainee picks an invoice uniformly at random and puts it back after a miss. The draws are independent, and the chance of success on each draw stays \GTenHLProbability{one-half}. There is no fixed limit on the number of draws. The trainee stops at the first qualifying invoice and includes that draw in the count.
& Verify that the probabilities for all possible draw counts are non-negative and sum to 1 over arbitrarily many draws. Find the chance that the first qualifying invoice takes \GTenHLValue{more than 3 draws}. \\
\end{longtable}
\subsubsection*{Complete worked solution}
\label{subsubsec:t2-c1-advanced-general-problem-an-infinite-model-solution}
\par\noindent\begin{minipage}{\linewidth}
\textbf{Step 1: Define and identify support.} $\GTenHLValue{T}$ is a \GTenHLValue{positive-integer trial count} with support $\GTenHLValue{\{1,2,\ldots\}}$. In the financial story, $\GTenHLValue{T}$ counts \GTenHLValue{calls}, including the successful call. A first sale on call $\GTenHLValue{t}$ requires $\GTenHLValue{t}-\GTenHLValue{1}$ failures followed by a success. Independence gives $p_{\GTenHLValue{T}}(\GTenHLValue{t})=(\GTenHLProbability{1/2})^{\GTenHLValue{t}-\GTenHLValue{1}}(\GTenHLProbability{1/2})=\GTenHLProbability{2^{-t}}$. This is the same infinite model; normalization must be checked.
\end{minipage}\par\medskip
\par\noindent\begin{minipage}{\linewidth}
\textbf{Step 2: Normalize.} The support has no last value, so first add the masses through a finite trial count $m$. Since $2^{-t}=2^{-(t-1)}-2^{-t}$, consecutive terms cancel in the sum.
\[
\begin{aligned}
\sum_{\GTenHLValue{t}=1}^{m}\GTenHLProbability{2^{-t}}
&=(1-\GTenHLProbability{2^{-1}})+(\GTenHLProbability{2^{-1}}-\GTenHLProbability{2^{-2}})+\cdots
+(\GTenHLProbability{2^{-(m-1)}}-\GTenHLProbability{2^{-m}})\\
&=1-\GTenHLProbability{2^{-m}}.
\end{aligned}
\]
As $m$ grows without bound, this partial sum approaches 1. Every $\GTenHLProbability{2^{-t}}$ is positive, so the infinite list defines a PMF. In the story, the total chance of a first sale on some finite call is 1.
\end{minipage}\par\medskip
\par\noindent\begin{minipage}{\linewidth}
\textbf{Step 3: Translate and calculate.} $P(\GTenHLValue{T}>\GTenHLValue{3})=1-\sum_{\GTenHLValue{t}=1}^{3}\GTenHLProbability{2^{-t}}=1-(\GTenHLProbability{1/2}+\GTenHLProbability{1/4}+\GTenHLProbability{1/8})=\GTenHLProbability{1/8}=\GTenHLProbability{0.1250}$.
\end{minipage}\par\medskip
\par\noindent\begin{minipage}{\linewidth}
\textbf{Step 4: Check and interpret.} The probability lies between 0 and 1. There is no last possible trial; infinite support does not mean infinite total probability. Financially, the chance of needing \GTenHLValue{more than 3 calls} for the first sale is \GTenHLProbability{12.5\%}. The waiting count is measured in \GTenHLValue{calls}; the probability has no units.
\end{minipage}\par\medskip
\noindent\textbf{Across the three narratives.} The variable $\GTenHLValue{T}$ counts \GTenHLValue{calls} through the first sale in Version 1, \GTenHLValue{quotes} through the first one within budget in Version 2, and \GTenHLValue{invoice draws} through the first qualifying invoice in Version 3. Each story gives a \GTenHLProbability{12.5\%} chance of needing \GTenHLValue{more than 3 attempts} to reach the first success.\par\medskip
% Graphs follow completed calculations; captions stay with their panels.
\noindent The graphs use Version 1. Write \GTenHLOutcome{F} for a failed call and \GTenHLOutcome{S} for a sale. The complete stopping outcomes are $\GTenHLOutcome{\mathrm{S}},\GTenHLOutcome{\mathrm{FS}},\GTenHLOutcome{\mathrm{FFS}},\ldots$; each ends at the first sale. The first two panels preview the first \GTenHLValue{five call counts}; further outcomes continue indefinitely.\par\medskip
\Needspace{6.8cm}
\par\noindent \begin{minipage}[t]{.485\linewidth}
\centering
\GTenOutcomeProbabilityGraph[width=\linewidth,height=4.7cm,rotate=35,ystep=.1,tail=1/32,xlabel={\GTenHLOutcome{Stopping outcome} ($\ldots$)},ylabel={\GTenHLProbability{Probability}},xhighlight=GTenHLOutcome,yhighlight=GTenHLProbability,color=GTenHLProbability!32]{S,FS,FFS,FFFS,FFFFS}{1/2,1/4,1/8,1/16,1/32}
\par\small \textbf{Outcome probabilities (Version 1, finite preview).} Later stopping outcomes have total \GTenHLProbability{probability $1/32$}. These bars sum to \GTenHLProbability{$31/32$}, without renormalization.
\end{minipage}\hfill
\begin{minipage}[t]{.485\linewidth}
\centering
\GTenOutcomeValueGraph[width=\linewidth,height=4.7cm,rotate=35,xlabel={\GTenHLOutcome{Stopping outcome} ($\ldots$)},ylabel={$\GTenHLValue{T}(\GTenHLOutcome{\omega})$ \GTenHLValue{(calls)}},xhighlight=GTenHLOutcome,yhighlight=GTenHLValue,color=GTenHLValue!32]{S,FS,FFS,FFFS,FFFFS}{1,2,3,4,5}
\par\small \textbf{Assigned call counts (Version 1).} Count every call, including the sale. Later outcomes map to \GTenHLValue{$6,7,\ldots$ calls}; their omitted probability is \GTenHLProbability{$1/32$}.
\end{minipage}\par\medskip

\Needspace{6.8cm}
\par\noindent \begin{minipage}[t]{.98\linewidth}
\centering
\GTenPMFGraph[width=\linewidth,height=4.7cm,xlabel={\GTenHLValue{$t$ (calls)}, continuing to $\GTenHLValue{9,10,\ldots}$},tail=1/256,ylabel={\GTenHLProbability{Probability}},xhighlight=GTenHLValue,yhighlight=GTenHLProbability,color=GTenHLProbability!32]{1,2,3,4,5,6,7,8}{1/2,1/4,1/8,1/16,1/32,1/64,1/128,1/256}
\par\small \textbf{Waiting-count PMF (Version 1, first eight values).} $p_{\GTenHLValue{T}}(\GTenHLValue{t})=\GTenHLProbability{2^{-t}}$ for $\GTenHLValue{t}=\GTenHLValue{1,2,\ldots}$; the support continues indefinitely to the right. The omitted probability is $P(\GTenHLValue{T}>\GTenHLValue{8})=\sum_{\GTenHLValue{t}=9}^{\infty}\GTenHLProbability{2^{-t}}=\GTenHLProbability{1/256}$. Displayed bars total $\GTenHLProbability{255/256}$ and are not renormalized; labels are rounded.
\end{minipage}\par\medskip

\secInicio[subsec:t2-c1-advanced-general-problem-an-infinite-model][sec:t2-c1-advanced]

\Needspace{8\baselineskip}
\subsection{Calls until a sale}
\label{subsec:t2-c1-calls-until-a-sale}
The \hyperref[subsec:t2-c1-advanced-general-problem-an-infinite-model]{preceding model} found $P(T>3)$. Here, use the same masses to find the complementary probability $P(T\leq3)$.

A sales team models $\GTenHLValue{T}$, the call number of its first sale, by $p_{\GTenHLValue{T}}(\GTenHLValue{t})=(\GTenHLProbability{1/2})^{\GTenHLValue{t}}$ for $\GTenHLValue{t}=\GTenHLValue{1,2,\ldots}$. It also records \GTenHLSecond{exact call duration} $\GTenHLSecond{D}$ in \GTenHLSecond{minutes}.

\Needspace{5\baselineskip}
\QuestionsTasks[sec:t2-c1-advanced]{
	{1}/{subsubsec:t2-c1-calls-until-a-sale-solution}/
	{Classify both variables; verify the waiting-count PMF and find $P(T\leq3)$.}
}
\subsubsection*{Complete worked solution}
\label{subsubsec:t2-c1-calls-until-a-sale-solution}

\textit{Markers here:} \GTenHLValue{$T$, call count}; \GTenHLSecond{$D$, duration in minutes}.

\par\noindent\begin{minipage}{\linewidth}
\textbf{Step 1: Define and distinguish.} $\GTenHLValue{T}$ counts calls and has countably infinite support. $\GTenHLSecond{D}$ measures time and is treated as continuous; its exact support contains intervals, rather than a list of isolated counts.
\end{minipage}\par\medskip
\par\noindent\begin{minipage}{\linewidth}
\textbf{Step 2: Check the supplied model.} All masses $\GTenHLProbability{2^{-t}}$ are positive. Their first $m$ terms sum to $1-\GTenHLProbability{2^{-m}}\to1$, so the infinite PMF is valid. It is compatible with independent calls with constant sale chance $\GTenHLProbability{1/2}$.
\end{minipage}\par\medskip
\par\noindent\begin{minipage}{\linewidth}
\textbf{Step 3: Calculate.} $P(\GTenHLValue{T}\leq\GTenHLValue{3})=\sum_{\GTenHLValue{t}=1}^{3}p_{\GTenHLValue{T}}(\GTenHLValue{t})=\GTenHLProbability{1/2}+\GTenHLProbability{1/4}+\GTenHLProbability{1/8}=\GTenHLProbability{7/8}=\GTenHLProbability{0.8750}$.
\end{minipage}\par\medskip
\par\noindent\begin{minipage}{\linewidth}
\textbf{Step 4: Interpret.} The first sale occurs \GTenHLValue{within three calls} with probability \GTenHLProbability{87.5\%}. There is no fixed maximum call number in the model; this does not make the count continuous.
\end{minipage}\par\medskip
% Graphs follow completed calculations; captions stay with their panels.
\noindent Only a distribution of numerical values is supplied, so the PMF below does not assume an underlying outcome list.\par\medskip
No distribution is supplied for the continuous duration $\GTenHLSecond{D}$; no PMF or outcome-to-duration graph can be inferred.\par\medskip
\Needspace{6.8cm}
\par\noindent \begin{minipage}[t]{.98\linewidth}
\centering
\GTenPMFGraph[width=\linewidth,height=4.7cm,xlabel={\GTenHLValue{$t$ (calls)}, continuing to $\GTenHLValue{9,10,\ldots}$},tail=1/256,ylabel={\GTenHLProbability{Probability}},xhighlight=GTenHLValue,yhighlight=GTenHLProbability,color=GTenHLProbability!32]{1,2,3,4,5,6,7,8}{1/2,1/4,1/8,1/16,1/32,1/64,1/128,1/256}
\par\small \textbf{Waiting-count PMF: first eight values only.} $p_{\GTenHLValue{T}}(\GTenHLValue{t})=\GTenHLProbability{2^{-t}}$ for $\GTenHLValue{t}=\GTenHLValue{1,2,\ldots}$; the support continues indefinitely to the right. The omitted probability is $P(\GTenHLValue{T}>\GTenHLValue{8})=\sum_{\GTenHLValue{t}=9}^{\infty}\GTenHLProbability{2^{-t}}=\GTenHLProbability{1/256}$. Displayed bars total $\GTenHLProbability{255/256}$ and are not renormalized; labels are rounded.
\end{minipage}\par\medskip

\secInicio[subsec:t2-c1-calls-until-a-sale][sec:t2-c1-advanced]

\Needspace{8\baselineskip}
\subsection{Delivery records: two numerical descriptions}
\label{subsec:t2-c1-delivery-records-two-numerical-descriptions}
Return to the idea from the \hyperref[subsec:t2-c1-independent-outcomes]{school sale}. One experiment can define several random variables.

A delivery has outcomes \GTenHLOutcome{early}, \GTenHLOutcome{on time}, \GTenHLOutcome{late} with probabilities $\GTenHLProbability{0.2},\GTenHLProbability{0.5},\GTenHLProbability{0.3}$. A bonus $\GTenHLValue{B}$ is \GTenHLValue{2 USD} for \GTenHLOutcome{early} or \GTenHLOutcome{on time} and \GTenHLValue{$-3$ USD} for \GTenHLOutcome{late}. An indicator $\GTenHLSecond{I}$ is \GTenHLSecond{1} if \GTenHLOutcome{late} and \GTenHLSecond{0} otherwise.

\QuestionsTasks[sec:t2-c1-advanced]{
	{1}/{subsubsec:t2-c1-delivery-records-two-numerical-descriptions-solution}/
	{Build both PMFs and contrast them with an exact delivery-time measurement.}
}
\subsubsection*{Complete worked solution}
\label{subsubsec:t2-c1-delivery-records-two-numerical-descriptions-solution}

\textit{Markers here:} \GTenHLValue{$B$, bonus in USD}; \GTenHLSecond{$I$, lateness indicator}.

\par\noindent\begin{minipage}{\linewidth}
\textbf{Step 1: Define and map.} $\GTenHLValue{B}$ measures \GTenHLValue{bonus in USD}, with support $\GTenHLValue{\{-3,2\}}$; $\GTenHLSecond{I}$ labels lateness, with support $\GTenHLSecond{\{0,1\}}$. \GTenHLOutcome{Early} and \GTenHLOutcome{on time} both map to $\GTenHLValue{B}=\GTenHLValue{2}$ and $\GTenHLSecond{I}=\GTenHLSecond{0}$. A \GTenHLOutcome{late} delivery maps to $\GTenHLValue{B}=\GTenHLValue{-3}$ and $\GTenHLSecond{I}=\GTenHLSecond{1}$. The indicator is unitless.
\end{minipage}\par\medskip
\par\noindent\begin{minipage}{\linewidth}
\textbf{Step 2: Check the categories.} The three exclusive outcomes cover all deliveries and their probabilities total 1. No equal-likelihood assumption is justified.
\end{minipage}\par\medskip
\par\noindent\begin{minipage}{\linewidth}
\textbf{Step 3: Calculate.} $p_{\GTenHLValue{B}}(\GTenHLValue{2})=\GTenHLProbability{0.2}+\GTenHLProbability{0.5}=\GTenHLProbability{0.7}$, $p_{\GTenHLValue{B}}(\GTenHLValue{-3})=\GTenHLProbability{0.3}$; $p_{\GTenHLSecond{I}}(\GTenHLSecond{0})=\GTenHLProbability{0.7}$, $p_{\GTenHLSecond{I}}(\GTenHLSecond{1})=\GTenHLProbability{0.3}$. Thus $P(\GTenHLValue{B}<\GTenHLValue{0})=P(\GTenHLSecond{I}=\GTenHLSecond{1})=\GTenHLProbability{0.3000}$.
\end{minipage}\par\medskip
\par\noindent\begin{minipage}{\linewidth}
\textbf{Step 4: Check and interpret.} Each PMF totals 1. The same uncertain delivery can generate different discrete variables with different units. Exact elapsed delivery time is a continuous measurement under an ideal model, even though its \GTenHLOutcome{early/on-time/late} classification is discrete. No distribution or exact time values are supplied, so no numerical mapping or PMF for elapsed time can be drawn.
\end{minipage}\par\medskip
% Graphs follow completed calculations; captions stay with their panels.
\Needspace{6.8cm}
\par\noindent \begin{minipage}[t]{.485\linewidth}
\centering
\GTenOutcomeProbabilityGraph[width=\linewidth,height=4.7cm,xlabel={\GTenHLOutcome{Outcome}},ylabel={\GTenHLProbability{Probability}},xhighlight=GTenHLOutcome,yhighlight=GTenHLProbability,color=GTenHLProbability!32]{early,on time,late}{.2,.5,.3}
\par\small \textbf{Delivery-outcome probabilities.} The bars refer to delivery categories.
\end{minipage}\hfill\begin{minipage}[t]{.485\linewidth}
\centering
\GTenOutcomeValueGraph[width=\linewidth,height=4.7cm,ylabel={$\GTenHLValue{B}(\GTenHLOutcome{\omega})$ \GTenHLValue{(USD)}},xlabel={\GTenHLOutcome{Outcome}},xhighlight=GTenHLOutcome,yhighlight=GTenHLValue,color=GTenHLValue!32]{early,on time,late}{2,2,-3}
\par\small \textbf{Assigned bonuses.} \GTenHLOutcome{Late} delivery gives \GTenHLValue{$-3$ USD}, shown below zero.
\end{minipage}\par\medskip
\Needspace{6.8cm}
\par\noindent \begin{minipage}[t]{.485\linewidth}
\centering
\GTenOutcomeValueGraph[width=\linewidth,height=4.7cm,ylabel={$\GTenHLSecond{I}(\GTenHLOutcome{\omega})$ \GTenHLSecond{(indicator)}},xlabel={\GTenHLOutcome{Outcome}},xhighlight=GTenHLOutcome,yhighlight=GTenHLSecond,color=GTenHLSecond!32]{early,on time,late}{0,0,1}
\par\small \textbf{Assigned lateness indicator.} Only \GTenHLOutcome{late} delivery maps to \GTenHLSecond{1}; the indicator has no units.
\end{minipage}\par\medskip
\Needspace{6.8cm}
\par\noindent \begin{minipage}[t]{.485\linewidth}
\centering
\GTenPMFGraph[width=\linewidth,height=4.7cm,xlabel={\GTenHLValue{$b$ (USD)}},ylabel={\GTenHLProbability{Probability}},xhighlight=GTenHLValue,yhighlight=GTenHLProbability,color=GTenHLProbability!32]{-3,2}{.3,.7}
\par\small \textbf{Bonus PMF.} \GTenHLOutcome{Early} and \GTenHLOutcome{on time} group at \GTenHLValue{2 USD}; a negative bonus still has a positive probability.
\end{minipage}\hfill\begin{minipage}[t]{.485\linewidth}
\centering
\GTenPMFGraph[width=\linewidth,height=4.7cm,xlabel={\GTenHLSecond{$i$ (indicator)}},ylabel={\GTenHLProbability{Probability}},xhighlight=GTenHLSecond,yhighlight=GTenHLProbability,color=GTenHLProbability!32]{0,1}{.7,.3}
\par\small \textbf{Indicator PMF.} \GTenHLOutcome{Early} and \GTenHLOutcome{on time} group at \GTenHLSecond{0}. Compare the outcome probabilities above.
\end{minipage}\par\medskip

\secInicio[subsec:t2-c1-delivery-records-two-numerical-descriptions][sec:t2-c1-advanced]

\end{document}
